% P16-C4: constructive hedge witnesses and domain-explicit completeness.
\section{Hedges, failure extraction and completeness}
\label{sec:p16-hedges-and-completeness}

Shpitser and Pearl's hedge argument supplies the obstruction behind
identification failure \cite[Definitions~5--6, Theorems~4--7 and
Corollaries~2--4, pp.~1222--1226]{ShpitserPearl2006JointIdentification}.
The following construction makes the countermodels, their support,
and their observed value sets explicit. In particular, a positive
countermodel on four values cannot silently settle a question whose
variables have prescribed binary value sets.

\begin{hbdefinition}{def:p16-hedge-identification-class}{The fixed-domain identification question}
Fix a mixed graph \(G\) on \(N\), finite nonempty observed value
sets \(S_i\), and the model class of
\cref{def:p16-latent-input-model}, allowing all finite independent
latent inputs with those declared edges. Functions may ignore declared
arguments. Put \(\mathbf S=(S_i)_{i\in N}\), and write
\(\mathcal M^+(G,\mathbf S)\) for its models whose observed
law is positive at every assignment in \(S_N\).

For disjoint \(X,Y\), the intervention kernel on \(Y\) under
\(X\) is \emph{identified in this class} if every two models in
the class with equal full observed laws have equal intervention
kernels. This quantifies over the model class and all its possible
observed laws. A counterexample at one positive law disproves this
claim; it need not disprove identification under additional restrictions
or for every particular observed law. A discrepancy on an event in
\(S_Y\) entails a discrepancy in at least one response cell.
\end{hbdefinition}

\begin{hbtheorem}{thm:p16-hedge-singleton-reduction}{proved=annex:p16.hedge-singleton-reduction}{Eliminating singleton observed domains}
Let \(C=\{i:|S_i|=1\}\), with unique assignment \(c\), and put
\(N_*=N\setminus C\), \(G_*=G[N_*]\), \(X_*=X\cap N_*\),
and \(Y_*=Y\cap N_*\). Substituting the constants in every
remaining equation gives a model on \(G_*\); it preserves the
observed law after deleting constant coordinates and every intervention
kernel after deleting constant action and response coordinates.
Conversely every model on \(G_*\) extends to one on \(G\)
with these properties. Positivity is preserved in both directions.
Consequently identification in \(\mathcal M^+(G,\mathbf S)\) is equivalent
to identification of the reduced kernel in the reduced class.
If \(Y_*\) is empty the requested response law is the unit mass.

Thus the algorithm used for the completeness statements below first
performs this reduction and then runs ID on the reduced positive law.
\end{hbtheorem}

\begin{hbdefinition}{def:p16-hedge-forests}{Confounded forests and hedges}
An \emph{edge subgraph} may delete vertices and either kind of edge;
it need not be an induced subgraph. A \(K\)-rooted
\emph{confounded forest}, or C-forest, is a nonempty bidirected-connected
mixed graph whose directed part is acyclic, whose vertices each have
at most one directed child, and whose vertices with no directed child
are exactly \(K\). These roots are directed sinks; they are distinct
from the parentless latent roots in a structural model.

For disjoint action and response sets \(X,Y\), a \emph{hedge}
is a pair \((F,F')\) of \(K\)-rooted C-forests, with vertex
sets \(T,S\), such that
\[
 F'\subseteq F\subseteq G,\qquad
 T\cap X\ne\varnothing,\qquad S\cap X=\varnothing,\qquad
 K\subseteq\operatorname{an}_{G[N\setminus X]}(Y).
\]
The last condition is the source's post-intervention ancestor condition:
an observed directed path from a nonaction vertex cannot enter a vertex
whose incoming arrows have been removed. Ancestry includes the target
itself. Both forests have exactly the same roots. In particular,
\(\varnothing\ne S\subsetneq T\).
\index{hedge!causal identification}\index{C-forest}
\end{hbdefinition}

\begin{hbtheorem}{thm:p16-hedge-failure-extraction}{proved=annex:p16.hedge-failure-extraction}{Extracting a hedge for the original query from every failure trace}
Suppose ID fails after singleton reduction. Use \(N,X_0,Y_0\)
for its reduced initial sets, and put
\(A=\operatorname{an}_{G[N\setminus X_0]}(Y_0)\).
At the recorded failed call, let \(R\) be its random set and
\(S\) its nonaction set. Exactly one of the following trace forms
applies:
\begin{enumerate}
\item There was no district split: \(S=A\), the local response is
      \(K=Y_0\), and \(A\) is one district.
\item There was one district split on this path: \(S=D\) for a
      maximal district \(D\) of \(G[A]\), and the local response
      is \(K=D\). All subsequent recursive edges use branches 2 or 7.
\end{enumerate}
In either case \(R\cap X_0\ne\varnothing\).
Choose one directed edge toward \(K\) at every vertex of
\(S\setminus K\), staying in \(S\), and one edge toward
\(S\) at every vertex of \(R\setminus S\). Stop at \(K\).
Choose a bidirected spanning tree on \(S\), and extend it to one
on \(R\). These selected edges give forests \(F,F'\), on
\(R,S\), which form a hedge for the \emph{original} query
\((X_0,Y_0)\). The saved trace therefore supplies the requisite
relationship to the original actions without claiming that all local
action sets are subsets of \(X_0\).
\end{hbtheorem}

\subsection{Strictly positive countermodels}

\begin{hbdefinition}{def:p16-hedge-parity-models}{The two forest models}
For a hedge with vertex sets \(T,S\) and roots \(K\), retain
a bidirected spanning tree on \(S\), extended to one on \(T\).
Associate an independent fair bit \(U_e\) with each selected
bidirected edge, as a latent input at its two endpoints. Choose
\(k_0\in K\) and an independent private bit \(E\) at \(k_0\),
with \(\mathbb P(E=1)=\epsilon\), \(0<\epsilon<1/2\).
All observed variables in these two models are bits. The symbol
\(\oplus\) denotes addition modulo two, with empty sum zero.

In model 1, every observed equation is the parity of its selected
directed parents, its incident edge bits, and also \(E\) at \(k_0\).
In model 2, vertices outside \(S\) keep these equations, while
vertices in \(S\) use only selected directed parents in \(S\)
and edge bits whose two endpoints lie in \(S\), again adding
\(E\) at \(k_0\). Other declared graph arguments are ignored.
\end{hbdefinition}

\begin{hbtheorem}{thm:p16-hedge-positive-parity}{proved=annex:p16.hedge-positive-parity}{Equal positive observations and different root parity}
The two forest models have the same strictly positive observed law
\[
 p(v_T)=2^{-(|T|-1)}
 \begin{cases}
  1-\epsilon,&\bigoplus_{k\in K}v_k=0,\\
  \epsilon,&\bigoplus_{k\in K}v_k=1.
 \end{cases}
\]
For every nonempty action set \(Z\subseteq T\setminus S\)
and every assignment \(z\), their root-parity probabilities differ:
\[
 \mathbb P^1_z\!\left(\bigoplus_{k\in K}V_k=1\right)=\tfrac12,
 \qquad
 \mathbb P^2_z\!\left(\bigoplus_{k\in K}V_k=1\right)=\epsilon.
\]
The construction uses only independent pair-specific latent bits
and one private bit. It is a permitted subclass of the larger
latent-input class, rather than a replacement theorem for its other models.
\end{hbtheorem}

\begin{hblemma}{thm:p16-hedge-alphabet-lift}{proved=annex:p16.hedge-alphabet-lift}{Lifting a positive witness to larger prescribed value sets}
Suppose two models use finite observed alphabets \(B_i\), have
equal strictly positive observed laws, and disagree on a specified
intervention response event. For each prescribed \(S_i\), suppose
there is a surjection \(d_i:S_i\to B_i\). Choose positive
weights \(w_i(s)\) summing to one on each fiber of \(d_i\).
There are models on the same declared graph and the prescribed
\(S_i\) whose common observed law is
\[
 \widetilde p(v)=p(d(v))\prod_i w_i(v_i)>0.
\]
Their decoded intervention law under \(x\) is the original
intervention law under \(d(x)\). Thus the discrepancy survives
on the inverse image of the original response event. In particular,
a positive binary witness lifts to every set of at least two values,
and a positive four-valued witness to every set of at least four values.
\end{hblemma}

\begin{hbdefinition}{def:p16-hedge-external-routing}{An external routing forest}
For a hedge \((F,F')\), an \emph{external routing forest} is a
selected subgraph of the directed part of \(G\) on \(K\cup H\),
where \(H\cap T=\varnothing\)
and \((K\cup H)\cap X=\varnothing\). Every vertex has at most
one selected child, every vertex is reached from \(K\), and
its sinks are a set \(J\subseteq Y\). No selected edge points
into \(K\). Every vertex of \(K\) therefore has exactly one
selected path to \(J\), allowing a zero-length path. Several
paths may merge. The graph with \(H=\varnothing\), no edges,
and \(J=K\) is allowed when \(K\subseteq Y\).
\end{hbdefinition}

\begin{hbtheorem}{thm:p16-hedge-binary-routing}{proved=annex:p16.hedge-binary-routing}{Transport of the positive binary discrepancy}
If a hedge has an external routing forest, its countermodels extend
to the whole declared graph with binary observed variables and
strictly positive equal observed laws, but different intervention
laws on \(Y\) under \(X\). For each \(h\in H\), use the
parity of its selected parents plus an independent private bit
\(E_h\), with \(0<\epsilon_h=\mathbb P(E_h=1)<1/2\).
The difference is witnessed by the parity on \(J\): its signed
expectation \(\mathbb E[(-1)^{\oplus_{j\in J}V_j}]\) is zero
in model 1 and
\[
 (1-2\epsilon)\prod_{h\in H}(1-2\epsilon_h)>0
\]
in model 2. This also proves nonidentification on any prescribed
value sets of size at least two by the alphabet lift.
\end{hbtheorem}

\begin{hbtheorem}{thm:p16-hedge-four-state-routing}{proved=annex:p16.hedge-four-state-routing}{A complete positive routing construction with four observed values}
Every hedge yields equal-observation, different-effect models in
\(\mathcal M^+(G,\mathbf S)\) whenever \(|S_i|\geq4\) for every
observed vertex. No external-routing condition is required.

For the four-valued construction write the variable at vertex \(i\)
as the pair \((C_i,L_i)\) of a core bit and a relay bit. The core
bits on \(T\) use the two forest models. Choose a directed forest
from all of \(K\) to \(Y\) inside \(G[N\setminus X]\).
Relay equations carry the parity of the core bits injected at \(K\)
to the chosen response sinks, with independent biased private noise
at every relay vertex. Core equations ignore all relay bits.
The annex supplies the total equations, equal-law and positivity
proofs, and the calculation after the original intervention.

More precisely, let \(H\) be the union of the selected routes,
including their initial roots and terminal responses. Four values
are required only at \(I=(T\setminus K)\cap H\); two values
suffice at every other vertex. At a root in \(K\), the core bit
can be computed inside its relay equation and omitted from its
observed output. The annex proves this coordinate elimination and
then applies the alphabet lift for these possibly different local
cardinality requirements.
\end{hbtheorem}

\subsection{Completeness with its exact domain}

\begin{hbtheorem}{thm:p16-hedge-completeness}{proved=annex:p16.hedge-completeness}{Failure and completeness on the proved fixed-domain classes}
First eliminate singleton domains. The following conclusions hold
for the remaining positive-law model class.
\begin{enumerate}
\item If all remaining value sets have at least four elements,
      ID succeeds exactly for the identified intervention kernels.
      Every failed trace constructs a positive counterexample for
      its original action and response sets.
\item If all remaining value sets have at least two elements and
      the response is the full nonaction set \(Y=N\setminus X\),
      the same completeness conclusion holds.
\item For arbitrary response sets and remaining value sets of size
      at least two, a failed trace whose extracted hedge has an
      external routing forest proves nonidentification of the
      original query. The same conclusion holds if a selected
      routing forest meets the more precise local cardinality
      requirements of \cref{thm:p16-hedge-four-state-routing}.
      Successful runs retain the previous soundness
      theorem on every finite nonempty value set.
\end{enumerate}
In the first two classes, do-calculus with ordinary probability
operations is complete for identified unconditional intervention
kernels: the successful ID trace gives the derivation.
For any individual failed query covered by the third clause, its
countermodels rule out a sound universal identification derivation.
\end{hbtheorem}

\begin{hbcorollary}{thm:p16-hedge-all-effects}{proved=annex:p16.hedge-all-effects}{Graphs identifying every intervention kernel}
After singleton elimination, every unconditional intervention kernel
is identified in the positive-law class if and only if no directed
edge has its endpoints in the same maximal district. This statement
holds for all remaining prescribed finite value sets of size at least
two. When the condition holds, every ID run succeeds and supplies
a do-calculus derivation. When it fails, a positive witness for a
full nonaction response kernel disproves identification of all effects.
\end{hbcorollary}

\begin{proofboundary}{def:p16-hedge-remaining-domain}{The remaining small-domain and zero-support obligations}
The routing construction retains both a core bit and a relay bit
at each nonroot forest vertex traversed by a selected route. It
does not prove the general fixed-binary or fixed-ternary case when
no available hedge and routing choice meets the proved local
cardinality requirements. Completing that case requires a countermodel
construction on those unchanged value sets, or a proved reduction
eliminating the extra channel at these remaining intersections while
preserving equal positive full observed laws and the original response
discrepancy.
Neither the alphabet lift nor the source's nonpositive parity law
supplies that reduction.

The source's general graphical completeness argument and its
do-calculus consequence motivate the constructions above; the
annex preserves the parity and failure-extraction proof routes
with these explicit domain refinements. An unrestricted execution
and completeness theorem at zero observed masses is also separate:
one must specify the recoverable action/support domain and justify
every conditional ratio. No structural intervention is made nonexistent
by a missing observational denominator.
\end{proofboundary}

\subsection{Calculated extraction, propagation and domain examples}

\begin{hbtheorem}{thm:p16-hedge-worked}{proved=annex:p16.hedge-worked}{A split failure, a positive propagated discrepancy, and a three-value lift}
Use graph \(A\to B,A\to C,B\to Z,C\to Z\), with
\(A\leftrightarrow B\leftrightarrow C\) and no other
bidirected edges. For \(P_a(z)\), ID splits the post-action
districts \(\{B,C\},\{Z\}\). The first child restricts to
\(\{A,B,C\}\) and fails. Select \(A\to B\), omitting
\(A\to C\); the extracted forests have vertex sets
\(\{A,B,C\},\{B,C\}\) and common roots \(\{B,C\}\).

Let \(U,V,E,E_Z\) be mutually independent bits, with \(U,V\)
fair and \(E,E_Z\) of one-probability \(1/4\). In model 1 set
\[
 A=U,\quad B=A\oplus U\oplus V\oplus E,\quad C=V,
 \quad Z=B\oplus C\oplus E_Z.
\]
Model 2 replaces the \(B\) equation by \(B=V\oplus E\).
Their common full observed law is given by the following table;
each listed mass applies to each of the four triples \((a,b,c)\)
with the stated parity \(q=b\oplus c\):
\[
\begin{array}{c|cc}
q&p(a,b,c,0)&p(a,b,c,1)\\\hline
0&9/64&3/64\\
1&1/64&3/64
\end{array}
\]
Both tables are strictly positive, but
\(\mathbb P^1_a(Z=1)=1/2\) and
\(\mathbb P^2_a(Z=1)=3/8\) for each \(a\).

For the positive bow witness \(A=U\), with model-1 response
\(Y=A\oplus U\oplus E\) and model-2 response \(Y=E\),
partition each three-element alphabet as \(\{0\},\{1,2\}\)
and weight the second fiber equally. The lifted common observed
matrix, with rows \(A=0,1,2\) and columns \(Y=0,1,2\), is
\[
 \frac1{32}\begin{pmatrix}12&2&2\\6&1&1\\6&1&1\end{pmatrix}.
\]
After any action on \(A\), the response event \(Y\in\{1,2\}\)
has probability \(1/2\) in model 1 and \(1/4\) in model 2.
If instead \(A\) has a singleton domain, its legal intervention
changes no equation value, and the bow response is identified by
its observed marginal after the proved elimination step.
\end{hbtheorem}
